\section{Identification of Causal Effects in CBNs}

This section investigates under which circumstances one can \emph{identify causal effects} and estimate them just from observational data alone under the (strong) assumption that the underlying causal graph is known. More generally, we ask the question when an interventional Markov kernel of a causal Bayesian network can be identified from the causal graph $G$ and the observational Markov kernel alone.

We will see that the main tool to allow for such statements is the global Markov property, see Theorem \ref{thm-gmp-mI-CBN}, applied to the causal Bayesian network that is augmented with further intervention variables.

We first study under which graphical conditions interventions don't have an effect or when one essentially can replace interventions with conditioning operations. These rules will be summarized as the \emph{three rules of do-calculus}. The main references are \cite{Pearl93as,Pearl93bs,Pearl09}, also see \cite{Pearl95s,FM20,For21}.

We then study under which graphical criteria one gets explicit \emph{adjustment formulas} to estimate interventional Markov kernels from observational ones. The literature mentions the \emph{backdoor criterion}, see \cite{Pearl93as,Pearl93bs,Pearl09}, 
the \emph{extended backdoor criterion}, see \cite{PearlPaz13,SWR10s}, the \emph{selection backdoor criterion}, see \cite{BTP14s},
criteria for \emph{selection without/partial external data}, see \cite{CB17s,CTB18s}, and all their generalizations to the cyclic case,  see \cite{FM20}, also see \cite{SP06s,PTKM15s,For21}.

Finally, we present the \emph{ID-algorithm}, which can decide just by processing the causal graph $G$ if an interventional Markov kernel can be identified by the observational one (under further assumptions, like strict positivity, etc.). If the algorithm does not output FAIL then it also presents a formula to estimate the queried interventional Markov kernel. 
The main references for the ID-algorithm are \cite{Pearl09,GP95s,Tian2002,TP02s,Tian04s,SP06,HV06s,HV08s,RERS17,FM20}.




\subsection{Do-Calculus}

\Joris{This whole section needs to be checked for the changed definition of intervention nodes}
\begin{Rem}[Recap]
    \label{rem-bcn-notation}
    Consider an IL-BCN:
    \[M= \lp G^+=\lp J,(V,U),E^+ \rp, \lp P_v ( X_v|\doit( X_{\Pa^{G^+}(v)})) \rp_{v \in V \cup U} \rp.  \]
        Then we get the joint  Markov kernel over all input, observed and unobserved output variables as follows:
        \[ P\lp X_V,X_U,X_J|\doit(X_J) \rp := \bigotimes_{v \in U \cup V} P_v\lp X_v|\doit(X_{\Pa^{G^+}(v)})\rp \otimes 
        \bigotimes_{j \in J} \delta(X_j|X_j).\]
    Further, for $D \ins J \cup V$ and $C \ins V \sm D$ we get the combined hard and soft interventions:
    \[ P\lp X_{V \sm D},X_U,X_{J \cup D},X_{I_C}| \doit(X_{I_C},X_{J \cup D}) \rp :=\] 
    \[    \bigotimes_{v \in V \sm (C \cup D)} P_v\lp X_v|\doit(X_{\Pa^{G^+}(v)})\rp \otimes 
    \bigotimes_{v \in C} P_v\lp X_v|\doit(X_{\Pa^{G^+}(v)},X_{I_v})\rp \otimes \]
     \[   \bigotimes_{v \in U} P_v\lp X_v|\doit(X_{\Pa^{G^+}(v)})\rp \otimes 
     \bigotimes_{j \in J \cup D} \delta(X_j|X_j) \otimes\bigotimes_{v \in C} \delta(X_{I_v}|X_{I_v})  ,\]
       where we need to reorder all the factors such that the product is in reverse order of a topological order and 
       where we use the following Markov kernels to model hard interventions as  soft interventions, $v\in C$: 
       \[ P_v\lp X_v|\doit\lp X_{\Pa^{G^+}(v)},X_{I_v}=x_{I_v}\rp\rp := 
         \left\{ \begin{array}{lll}
                 P_v\lp X_v|\doit\lp X_{\Pa^{G^+}(v)}\rp\rp, & \text{ if } & x_{I_v} = \star, \\
    \delta(X_v|X_v =x_{I_v}),  & \text{ if } & x_{I_v} \neq \star. 
                \end{array}\right.\]
  Finally we can also marginalize (i.e.\ integrating out) and condition to get: 
  \[ P\lp X_A|X_B, \doit(X_{J \cup D},X_{I_C}) \rp ,\]
for any $A,B,C,D \ins J \cup V$.\\
For more suggestive formulas later on we also freely permute the order of symbols behind the conditioning line, e.g.:
\[ P\lp X_A|\doit(X_F),X_B, \doit(X_D) \rp := P\lp X_A|X_B, \doit(X_D,X_F) \rp.\]
Please note that no matter in which order we write the $\doit$-part and conditioning part behind the conditioning line $|$, we always assume that we perform the intervention ($\doit$) first and afterwards condition.\\
We will also make use of the following CADMG:
\[ G_{\doit(I_C,D)} = (G^+_{\doit(I_C,D)})^{\sm U}.  \]
W.l.o.g.\ we can assume: $C \cap D = \emptyset$.
\end{Rem}

\begin{Lem}
    \label{lem:int-null-set}
    For pairwise disjoint $B,C \ins V$ and $D \ins V \cup J$ and a measurable subset $N \ins \Xcal_{B \cup C\cup D \cup J}$ the following statements are equivalent:
    \begin{enumerate}
        \item $N$ is a $P(X_B,X_C|\doit(X_{I_B},X_{D\cup J}))$-null set.
        \item For every decomposition $B=B_1 \dcup B_2$ the set $N$ is a 
            $P(X_B,X_C|\doit(X_{B_2},X_{D\cup J}))$-null set.
        \item For every decomposition $B=B_1 \dcup B_2$ the set $N$ is a 
            $P(X_{B_1},X_C|\doit(X_{B_2},X_{D\cup J}))$-null set.
    \end{enumerate}
\begin{proof}
    Every value $x_{I_B}=(x_{I_v})_{v \in B} \in \Xcal_{I_B}$ defines a decomposition $B=B_1 \dcup B_2$ via:
    \[ B_1 := \lC v \in B \st x_{I_v} = \star \rC, \qquad B_2 := \lC v \in B \st x_{I_v} \in \Xcal_v \rC.  \]
    So running through all values $x_{I_B} \in \Xcal_{I_B}$ is the same as running through all subsets $B_2 \ins B$ and all values $x_{B_2} \in \Xcal_{B_2}$, while putting $x_{I_{B_1}}=\star$ for $B_1 = B \sm B_2$.
    Furthermore, we have the following identities:
    \begin{align*}
          &   P((X_B,X_C) \in N_{x_{D\cup J}}|\doit(X_{I_B}=x_{I_B},X_{D\cup J}=x_{D\cup J})) \\
        &=   P((X_B,X_C) \in N_{x_{D\cup J}}|\doit(X_{I_{B_1}} = \star ,X_{I_{B_2}}=x_{B_2},X_{D\cup J}=x_{D\cup J})) \\
     &=   P((X_B,X_C) \in N_{x_{D\cup J}}|\doit(X_{B_2}=x_{B_2},X_{D\cup J}=x_{D\cup J})) \\
     &= \lp P(X_{B_1},X_C|\doit(X_{B_2}=x_{B_2},X_{D\cup J}=x_{D\cup J})) \otimes \delta(X_{B_2}|X_{B_2}=x_{B_2})\rp\lp N_{x_{D\cup J}} \rp \\
     &= P((X_{B_1},X_C) \in N_{(x_{B_2},x_{D\cup J})}|\doit(X_{B_2}=x_{B_2},X_{D\cup J}=x_{D\cup J})).
    \end{align*}
    So the first line vanishes for all values $x_{I_B} \in \Xcal_{I_B}$ and $x_{D\cup J} \in \Xcal_{D\cup J}$ if and only if any other line vanishes
    for all subsets $B_2 \ins B$ and all values $x_{B_2} \in \Xcal_{B_2}$ and $x_{D\cup J} \in \Xcal_{D\cup J}$.
    This shows the claim.
\end{proof}
\end{Lem}


\begin{Cor}[Do-calculus]\label{cor:do-calculus}
    Consider an IL-BCN:
    \[M= \lp G^+=\lp J,(V,U),E^+ \rp, \lp P_v ( X_v|\doit( X_{\Pa^{G^+}(v)})) \rp_{v \in V \cup U} \rp.\]
    Let $A,B,C \ins V$ and $D \ins J \cup V$ be such that $A,B,C,D$ are pairwise disjoint.\\
    \Joris{In rule 3, should we allow for $B \subseteq V \cup W \cup J$?}
    Then we have the following 3 rules relating marginal conditional to marginal interventional Markov kernels:
    \begin{enumerate}
        \item \underline{Insertion/deletion of observation}: If we have:
            \[ A \Perp^d_{G_{\doit(D)}} B \given C \cup D,   \]
            then there exists a Markov kernel:
            \[ P\lp X_A|\cancel{X_B},X_C,\doit(X_{D\cup \cancel{J}})\rp:\;  \Xcal_C \times \Xcal_{D} \dshto \Xcal_A, \]
            that is a version of:
            \[P\lp X_A|X_{B_2},X_C,\doit(X_{D\cup J})\rp,\] 
            for every subset $B_2 \ins B$ simultaneously.
            Note that this Markov kernel is only dependent on $x_C$ and $x_D$, and constant in $x_{J\setminus D}$.
        \item[] Such a Markov kernel is unique up to a measurable $P(X_C|\doit(X_{D\cup J}))$-null set in $\Xcal_{C\cup D}$. 
         \item \underline{Action/observation exchange}: If we have:
            \[ A \Perp^d_{G_{\doit(I_B,D)}} I_B \given B \cup C \cup D,   \]
            then there exists a Markov kernel:
            \[P\lp X_A|\cancel{\doit}(X_B),X_C,\doit(X_{D\cup \cancel{J}})\rp: \;\Xcal_B \times \Xcal_C \times \Xcal_{D} \dshto \Xcal_A,\]
            that is a version of:            
            \[ P\lp X_A|X_{B_1},\doit(X_{B_2}),X_C,\doit(X_{D\cup J})\rp, \]
            for every decomposition: $B=B_1\dcup B_2$, simultaneously.
        \item[] Such a Markov kernel is unique up to a measurable $P(X_B,X_C|\doit(X_{I_B},X_{D\cup J}))$-null set in $N \ins \Xcal_{B \cup C\cup D}$, i.e.\ $N$ is a $P(X_{B_1},X_C|\doit(X_{B_2},X_{D\cup J}))$-null set
            for every decomposition $B=B_1\dcup B_2$ simultaneously.
        \item \underline{Insertion/deletion of action}: If we have:
            \[ A \Perp^d_{G_{\doit(I_B,D)}} I_B \given C  \cup D,   \]
            then there exists a Markov kernel: 
            \[P\lp X_A|\cancel{\doit(X_B)},X_C,\doit(X_{D\cup \cancel{J}})\rp:\; \Xcal_C \times \Xcal_{D} \dshto \Xcal_A,\] 
            that is a version of:  
            \[ P\lp X_A|\doit(X_{B_2}),X_C,\doit(X_{D\cup J})\rp\]
            for every subset $B_2 \ins B$ simultaneously. 
            Note that this Markov kernel is only dependent on $x_C$ and $x_D$, and constant in $x_{J \sm D}$.
        \item[] Such a Markov kernel is unique up to a measurable $P\lp X_C|\doit(X_{I_B},X_{D\cup J})\rp$-null set $N \ins \Xcal_{C\cup D}$, i.e.\ $N$ is a $P(X_C|\doit(X_{B_2},X_{D\cup J}))$-null set for every subset $B_2 \ins B$ simultaneously.
    \end{enumerate}
\begin{proof}
    We make use of the global Markov property (GMP), theorem \ref{thm-gmp-mI-CBN}.

1.) The assumption:
            \[ A \Perp^d_{G_{\doit(D)}} B \given C \cup D,   \]
implies the conditional independence by GMP \ref{thm-gmp-mI-CBN}:
  \[ X_A \Indep_{P(X_V|\doit(X_{D\cup J}))} X_B \given X_C,X_D.\]
So  we get the following factorization, where we can omit the deterministic variables from $X_D$ on the left of the conditioning lines:
  \[ P\lp X_A,X_B,X_C|\doit(X_{D\cup J})\rp = Q(X_A|X_C,X_D) \otimes P\lp X_B,X_C|\doit(X_{D\cup J})\rp,\]
for some Markov kernel $Q(X_A|X_C,X_D)$. Here
$Q(X_A|X_C,X_D)$ serves as a version of the conditional Markov kernel:
  \[P\lp X_A|X_B,X_C,\doit(X_{D\cup J})\rp.\]
If we marginalize out $X_{B_1}$ for any decomposition $B=B_1 \dcup B_2$ 
in the above factorization we also get:
  \[ P\lp X_A,X_C,X_{B_2}|\doit(X_{D\cup J})\rp = Q(X_A|X_C,X_D) \otimes P\lp X_C,X_{B_2}|\doit(X_{D\cup J})\rp,\]
showing that $Q(X_A|X_C,X_D)$ is also a version of:
  \[P\lp X_A|X_{B_2},X_C,\doit(X_{D\cup J})\rp.\]
In particular, this holds for $B_2 = \emptyset$. This shows all the claimed properties for $Q(X_A|X_C,X_D)$.

Now consider another Markov kernel $K(X_A|X_C,X_D)$ and the measurable sets:
\begin{align*}
    \tilde N &:= \lC x_{C \cup D} \in \Xcal_{C \cup D} \st Q(X_A|X_C=x_C,X_D=x_D) \neq K(X_A|X_C=x_C,X_D=x_D)  \rC,\\
     N &:= \lC x_{C \cup D \cup J} \in \Xcal_{C \cup D \cup J} \st Q(X_A|X_C=x_C,X_D=x_D) \neq K(X_A|X_C=x_C,X_D=x_D)  \rC \\
       &= \tilde N \times \Xcal_{J \sm D}.
\end{align*}
If $K(X_A|X_C,X_D)$ is a version of: 
\[P(X_A|X_{B_2},X_C,\doit(X_{D\cup J})),\] 
for every subset $B_2 \ins B$ simultaneously, then this holds, in particular, for $B_2 = \emptyset$.
Since conditional Markov kernels are essentially unique, by Theorem \ref{thm-regular-conditional-Markov-kernel}, we have that $N$ is a $P(X_C|\doit(X_{D\cup J}))$-null set.
Note that then automatically the product $\Xcal_{B_2} \times N$ is a $P(X_{B_2},X_C|\doit(X_{D\cup J}))$-null sets for every $B_2 \ins B$. \\


2.) The assumption:
            \[ A \Perp^d_{G_{\doit(I_B,D)}} I_B \given B \cup C \cup D,   \]
implies the conditional independence by GMP \ref{thm-gmp-mI-CBN}:
  \[ X_A \Indep_{P(X_V|\doit(X_{I_B},X_{D\cup J}))} X_{I_B} \given X_B,X_C,X_D.\]
So we have the following factorization:
\begin{align}
    P\lp X_A,X_B,X_C|\doit(X_{I_B},X_{D\cup J})\rp &= Q(X_A|X_B,X_C,X_D) \otimes P\lp X_B,X_C|\doit(X_{I_B},X_{D\cup J})\rp,\label{eq:2nd-do-int-1}  
\end{align}
for some Markov kernel $Q(X_A|X_B,X_C,X_D)$, which serves as a version of the conditional Markov kernel:
  \[P\lp X_A|X_B,X_C,\doit(X_{I_B},X_{D\cup J})\rp,\]
and which is independent of $X_{I_B}$.

We fist claim that for a Markov kernel $Q(X_A|X_B,X_C,X_D)$ the equation \ref{eq:2nd-do-int-1} is equivalent to the system of equations \ref{eq:2nd-do-int-2} indexed by subsets $B_2 \ins B$ and with $B_1 := B \sm B_2$:
\begin{align}
    P\lp X_A,X_{B_1},X_C|\doit(X_{B_2},X_{D\cup J})\rp &= Q(X_A|X_B,X_C,X_D) \otimes P\lp X_{B_1},X_C|\doit(X_{B_2},X_{D\cup J})\rp. \label{eq:2nd-do-int-2} 
\end{align}
Indeed, we can look at the different input values for
$X_{I_{B}} = (X_{I_{B_1}},X_{I_{B_2}})$ in equation \ref{eq:2nd-do-int-1}.
For $B_1$ we put: $X_{I_{B_1}}= \star = (\star)_{v \in B_1}$ and for $B_2$ we take values: 
$X_{I_{B_2}}=x_{B_2} \in \Xcal_{B_2}$.
This implies: 
\begin{align*} 
  & P\lp X_A,X_B,X_C|\doit(X_{{B_2}}=x_{B_2},X_{D\cup J})\rp \\
  &= P\lp X_A,X_B,X_C|\doit(X_{I_{B_1}}=\star,X_{I_{B_2}}=x_{B_2},X_{D\cup J})\rp \\
  &= Q(X_A|X_B,X_C,X_D) \otimes P\lp X_B,X_C|\doit(X_{I_{B_1}}=\star,X_{I_{B_2}}=x_{B_2},X_{D\cup J})\rp,\\
  &= Q(X_A|X_B,X_C,X_D) \otimes P\lp X_B,X_C|\doit(X_{{B_2}}=x_{B_2},X_{D\cup J})\rp.
\end{align*}
So we get the equations:
\begin{align*}
    P\lp X_A,X_B,X_C|\doit(X_{B_2},X_{D\cup J})\rp 
  &= Q(X_A|X_B,X_C,X_D) \otimes P\lp X_B,X_C|\doit(X_{B_2},X_{D\cup J})\rp, 
\end{align*}
  where we can further marginalize out the deterministic $X_{B_2}$:
  \begin{align*} P\lp X_A,X_{B_1},X_C|\doit(X_{B_2},X_{D\cup J})\rp 
  &= Q(X_A|X_B,X_C,X_D) \otimes P\lp X_{B_1},X_C|\doit(X_{B_2},X_{D\cup J})\rp. %\label{eq:2nd-do-int-3} 
  \end{align*}
  Note  that we can also go back by multiplying with $\delta(X_{B_2}|X_{B_2})$. This shows the intermediate claim.

The equation \ref{eq:2nd-do-int-2} already implies that $Q(X_A|X_B,X_C,X_D)$ is a version of the conditional Markov kernel:
\begin{align} P\lp X_A|X_{B_1},X_C,\doit(X_{B_2},X_{D\cup J})\rp, \label{eq:2nd-do-int-4} 
\end{align}
for every decomposition: $B=B_1 \dcup B_2$ simultaneously.

Now consider another Markov kernel $K(X_A|X_B,X_C,X_D)$ and the measurable sets:
\begin{align*}
    \tilde N &:= \bigl\{ x_{B \cup C \cup D} \in \Xcal_{B \cup C \cup D} \,\big|\,  Q(X_A|X_B=x_B,X_C=x_C,X_D=x_D) \\ 
    & \qquad \qquad  \neq K(X_A|X_B=x_B,X_C=x_C,X_D=x_D)  \bigr\},\\
  N &:= \tilde N \times \Xcal_{J \sm D}.
\end{align*}
If $K(X_A|X_B,X_C,X_D)$ now is also a version of the conditional Markov kernel \ref{eq:2nd-do-int-4} 
for every decomposition $B=B_1 \dcup B_2$ simultaneously, then $N$ is a 
$P(X_{B_1},X_C|\doit(X_{B_2},X_{D\cup J}))$-null set for every decomposition $B=B_1 \dcup B_2$, 
because conditional Markov kernels are essentially unique, see Theorem \ref{thm-regular-conditional-Markov-kernel}. 
By Lemma \ref{lem:int-null-set} this statement is equivalent for $N$ to be a 
$P(X_B,X_C|\doit(X_{I_B},X_{D\cup J}))$-null set.\\


3.) The assumption:
            \[ A \Perp^d_{G_{\doit(I_B,D)}} I_B \given C  \cup D,   \]
implies the conditional independence by GMP \ref{thm-gmp-mI-CBN}:
\[ X_A \Indep_{P(X_V|\doit(X_{I_B},X_{D\cup J}))} X_{I_B} \given X_C,X_D.\]
So we have the following factorization:
\begin{align} 
    P\lp X_A,X_C|\doit(X_{I_B},X_{D\cup J})\rp &= Q(X_A|X_C,X_D) \otimes P\lp X_C|\doit(X_{I_B},X_{D\cup J})\rp,
\label{eq:3nd-do-int-1} 
\end{align}
for some Markov kernel $Q(X_A|X_C,X_D)$, which serves as a version of the conditional Markov kernel:
\[P\lp X_A|X_C,\doit(X_{I_B},X_{D\cup J})\rp,\]
and which is independent of $X_{I_B}$.\\ 
We can now look at the different input values for any decomposition: $B=B_1 \dcup B_2$.
For this we put: $X_{I_{B_1}}= \star = (\star)_{v \in B_1}$ and  $ X_{I_{B_2}}=x_{B_2} \in \Xcal_{B_2}$.
This implies:
\begin{align*} 
  & P\lp X_A,X_C|\doit(X_{{B_2}}=x_{B_2},X_{D\cup J})\rp \\
  &= P\lp X_A,X_C|\doit(X_{I_{B_1}}=\star,X_{I_{B_2}}=x_{B_2},X_{D\cup J})\rp \\
  &= Q(X_A|X_C,X_D) \otimes P\lp X_C|\doit(X_{I_{B_1}}=\star,X_{I_{B_2}}=x_{B_2},X_{D\cup J})\rp,\\
  &= Q(X_A|X_C,X_D) \otimes P\lp X_C|\doit(X_{{B_2}}=x_{B_2},X_{D\cup J})\rp.
\end{align*}
So we get for every subset $B_2 \ins B$:
\begin{align} 
    P\lp X_A,X_C|\doit(X_{B_2},X_{D\cup J})\rp &= Q(X_A|X_C,X_D) \otimes P\lp X_C|\doit(X_{B_2},X_{D\cup J})\rp, \label{eq:3nd-do-int-2} 
\end{align}
which shows that $Q(X_A|X_C,X_D)$ is a version of the conditional Markov kernel:
\begin{align} P\lp X_A|X_C,\doit(X_{B_2},X_{D\cup J})\rp, \label{eq:3nd-do-int-3}  \end{align}
for every subset $B_2 \ins B$ simultaneously.

Now consider another Markov kernel $K(X_A|X_C,X_D)$ and the measurable sets:
\begin{align*}
    N &:= \lC x_{C \cup D} \in \Xcal_{C \cup D} \st Q(X_A|X_C=x_C,X_D=x_D) \neq K(X_A|X_C=x_C,X_D=x_D)  \rC,\\
    N_{B_2} &:= \Xcal_{B_2} \times \tilde N \times \Xcal_{J \sm D}.
\end{align*}
Now assume that $K(X_A|X_C,X_D)$ is a version of the conditional Markov kernel \ref{eq:3nd-do-int-3} 
for every subset $B_2 \ins B$ simultaneously. 
Then for every subset $B_2 \ins B$ set $N_{B_2}$ is a
$P(X_C|\doit(X_{B_2},X_{D\cup J}))$-null set, because of the essential uniqueness of conditional Markov kernels, see Theorem \ref{thm-regular-conditional-Markov-kernel}. More concretely, for $x_{B_2} \in \Xcal_{B_2}$ 
and $x_{D \cup J} \in \Xcal_{D \cup J}$ we get the equations:
\begin{align*}
    0 &=P(X_C \in (N_{B_2})_{(x_{B_2},x_D)} |\doit(X_{B_2}=x_{B_2},X_{D\cup J}=x_{D \cup J})) \\
      &=P(X_C \in N_{x_D} |\doit(X_{B_2}=x_{B_2},X_{D\cup J}=x_{D \cup J}))\\
      &=P(X_C \in N_{x_D} |\doit(X_{I_{B_1}}=\star, X_{I_{B_2}}=x_{B_2},X_{D\cup J}=x_{D \cup J}))\\
      &=P(X_C \in N_{x_D} |\doit(X_{I_B}=(\star,x_{B_2}),X_{D\cup J}=x_{D \cup J})).
\end{align*}
Since we have this for all decompositions $B=B_1 \dcup B_2$ and all values $x_{B_2} \in \Xcal_{B_2}$ 
we are running through all values $x_{I_B} \in \Xcal_{I_B}$. 
This shows that $N$ is a $P(X_C|\doit(X_{I_B},X_{D\cup J}))$-null set in $\Xcal_{C \cup D}$.
\end{proof}
\end{Cor}

%\Joris{Add remark about completeness: although the do-calculus has been shown to be complete for CBNs without input variables, it is not complete for simple SCMs with input variables.}


\begin{Cor}[Almost-sure do-calculus]
    \label{cor:as-do-calculus}
    Consider the situations from Corollary \ref{cor:do-calculus}.
    Further assume that we have reference measures $\mu_v$ on $\Xcal_v$ for every $v \in V$ that are each equivalent to probability measure (in terms of absolute continuity). 
    We then put $\mu_F := \bigotimes_{v \in F} \mu_v$ for $F \ins V$.
\begin{enumerate}
\item \underline{Insertion/deletion of observation}: Assume:
        \[ A \Perp^d_{G_{\doit(D)}} B \given C \cup D. \]
        For a fixed finite index set $I$ consider subsets $B^{(i)} \ins B$, for $i \in I$, and  pick an arbitrary version of a conditional Markov kernel:
            \[P\lp X_A|X_{B^{(i)}},X_C,\doit(X_{D\cup J})\rp:\;\Xcal_{B \cup C \cup D \cup J} \to \Xcal_{B^{(i)} \cup C \cup D \cup J} \to \Prob(\Xcal_A), \]
            of $P(X_A,X_{B^{(i)}},X_C|\doit(X_{D\cup J}))$.
            Then there exists a measurable $P(X_B,X_C|\doit(X_{D\cup J}))$-null set
        $N \ins \Xcal_{B \cup C \cup D \cup J}$,
        such that all those Markov kernels are equal on the complement $N^\cmpl$.
    \item[] Note that if $\mu_{B \cup C} \ll P(X_B,X_C|\doit(X_{D\cup J}))$ then $N$ is also a $\mu_{B \cup C}$-null set, i.e.\ for every $x_{D \cup J} \in \Xcal_{D \cup J}$ we have:
        $\mu_{B \cup C} (N_{x_{D \cup J}}) =0.$
    \item[] If we also have the reverse $ P(X_B,X_C|\doit(X_{D\cup J})) \ll \mu_{B \cup C}$ then changing 
        the above conditional Markov kernels on a $\mu_{B \cup C}$-null set $N$ maintains their corresponding version of conditional Markov kernel.\footnote{Note that the absolute continuities: $\mu_{B \cup C}\ll P(X_B,X_C|\doit(X_{D\cup J})) \ll \mu_{B \cup C}$ hold if $P(X_B,X_C|\doit(X_{D\cup J}))$ has a strictly positive Doob-Radon-Nikodym derivative w.r.t.\ $\mu_{B \cup C}$.}

\item \underline{Action/observation exchange}: Assume:
            \[ A \Perp^d_{G_{\doit(I_B,D)}} I_B \given B \cup C \cup D.   \]
            For a fixed finite index set $I$ consider decompositions
            $B=B_1^{(i)} \dcup B_2^{(i)}$, for $i \in I$,
            and pick an arbitrary version of a conditional Markov kernel:
            \[ P(X_A|X_{B_1^{(i)}},\doit(X_{B_2^{(i)}}),X_C,\doit(X_{D\cup J})): 
            \;\Xcal_{B \cup C \cup D \cup J} \to \Prob(\Xcal_A), \]
            of $P(X_A,X_{B_1^{(i)}},X_C|\doit(X_{B_2^{(i)}},X_{D\cup J})) \otimes \mu_{B_2^{(i)}}$ 
            and assume the following absolute continuity:
            \[ \mu_{B \cup C} \ll  P(X_{B_1^{(i)}},X_C|\doit(X_{B_2^{(i)}},X_{D\cup J})) \otimes \mu_{B_2^{(i)}}. \]
        Then there exists a measurable $\mu_{B \cup C}$-null set
        $N \ins \Xcal_{B \cup C \cup D \cup J}$,
        such that all those conditional Markov kernels are equal on the complement $N^\cmpl$.
    \item[] If we also assume the reverse absolute continuity:
         \[  P(X_{B_1^{(i)}},X_C|\doit(X_{B_2^{(i)}},X_{D\cup J})) \otimes \mu_{B_2^{(i)}} \ll \mu_{B \cup C}, \]
         then all those conditional Markov kernels are versions of each other.\footnote{Note that the absolute continuities: $ \mu_{B \cup C} \ll P(X_{B_1^{(i)}},X_C|\doit(X_{B_2^{(i)}},X_{D\cup J})) \otimes \mu_{B_2^{(i)}} \ll \mu_{B \cup C} $ hold if the absolute continuities: $\mu_{B_1^{(i)} \cup C} \ll P(X_{B_1^{(i)}},X_C|\doit(X_{B_2^{(i)}},X_{D\cup J})) \ll \mu_{B_1^{(i)} \cup C} $ hold, which hold if $P(X_{B_1^{(i)}},X_C|\doit(X_{B_2^{(i)}},X_{D\cup J}))$ has a strictly positive Doob-Radon-Nikodym derivative w.r.t.\ $\mu_{B_1^{(i)} \cup C}$.}

\item \underline{Insertion/deletion of action}: Assume:
            \[ A \Perp^d_{G_{\doit(I_B,D)}} I_B \given C  \cup D.  \]
            For a fixed finite index set $I$ consider subsets $B^{(i)} \ins B$, for $i \in I$, and  pick an arbitrary version of a conditional Markov kernel:
            \[ P(X_A|\doit(X_{B^{(i)}}),X_C,\doit(X_{D\cup J})):\; \Xcal_{B \cup C \cup D \cup J} \to \Xcal_{B^{(i)} \cup C \cup D \cup J} \to \Prob(\Xcal_A), \]
            of $P(X_A,X_C|\doit(X_{B^{(i)}},X_{D\cup J}))$ and assume the following absolute continuity:
            \[ \mu_C \ll  P(X_C|\doit(X_{B^{(i)}},X_{D\cup J})). \]
            Then there exists a measurable $\mu_C$-null set
        $N \ins \Xcal_{B \cup C \cup D \cup J}$,
        such that all those conditional Markov kernels are equal on the complement $N^\cmpl$.
    \item[] If we also assume the reverse absolute continuity:
        \[ P(X_C|\doit(X_{B^{(i)}},X_{D\cup J})) \ll \mu_C, \]
        then all those conditional Markov kernels are versions of each other.\footnote{Note that absolute continuities: $\mu_C \ll P(X_C|\doit(X_{B^{(i)}},X_{D\cup J})) \ll \mu_C$ hold if $P(X_C|\doit(X_{B^{(i)}},X_{D\cup J}))$ has a strictly positive Doob-Radon-Nikodym derivative w.r.t.\ $\mu_C$.}
\end{enumerate}
\begin{proof}
 W.l.o.g.\ we can assume all $\mu_v$ to be probability measures.

 1.) Let $Q(X_A|X_C,X_D)$ be the Markov kernel from Corollary \ref{cor:do-calculus} point 1.
Recall that conditional Markov kernels are essentially unique by Theorem \ref{thm-regular-conditional-Markov-kernel}. 
    This shows that the set:
\begin{align*}
    \tilde N^{(i)} := \bigl\{ x_{B^{(i)} \cup C \cup D \cup J} \in \Xcal_{B^{(i)} \cup C \cup D \cup J} \,\big|\,
       & Q(X_A|X_C=x_C,X_D=x_D)\\ 
       &\neq P(X_A|X_{B^{(i)}}=x_{B^{(i)}},X_C=x_C,\doit(X_{D\cup J}=x_{D \cup J})) \bigr\},
\end{align*}
is a (measurable) $P(X_{B^{(i)}},X_C|\doit(X_{D\cup J}))$-null set. So the lifted set:
\begin{align*}
    N^{(i)} := \bigl\{ x_{B \cup C \cup D \cup J} \in \Xcal_{B \cup C \cup D \cup J} \,\big|\,
       & Q(X_A|X_C=x_C,X_D=x_D)\\ 
       &\neq P(X_A|X_{B^{(i)}}=x_{B^{(i)}},X_C=x_C,\doit(X_{D\cup J}=x_{D \cup J})) \bigr\},
\end{align*}
is then a (measurable) $P(X_B,X_C|\doit(X_{D\cup J}))$-null set. Then also the finite union:
 \[ N := \bigcup_{i \in I} N^{(i)}  \ins \Xcal_{B \cup C \cup D \cup J},  \]
 is a (measurable) $P(X_B,X_C|\doit(X_{D\cup J}))$-null set as well. 
 Note that on the complement $N^\cmpl$ all Markov kernels agree with $Q(X_A|X_C,X_D)$ and are thus all equal on $N^\cmpl$.\\

2.) Consider the Markov kernel $Q(X_A|X_B,X_C,X_D)$ from Corollary \ref{cor:do-calculus} point 2
    and for $i \in I$ the measurable set:
\begin{align*}
     N^{(i)} :=  \bigl\{ & x_{B \cup C \cup D \cup J} \in \Xcal_{B \cup C \cup D \cup J} \,\big|\,
        Q(X_A|X_B=x_B,X_C=x_C,X_D=x_D)\\ 
       &\neq P(X_A|X_{B_1^{(i)}}=x_{B_1^{(i)}},\doit(X_{B_2^{(i)}}=x_{B_2^{(i)}}),X_C=x_C,\doit(X_{D\cup J}=x_{D\cup J}))  \bigr\}.
\end{align*}
Again, by the essential uniqueness of conditional Markov kernels, Theorem \ref{thm-regular-conditional-Markov-kernel}, the set $N^{(i)}$ is 
%a $P(X_{B_1^{(i)}},X_C|\doit(X_{B_2^{(i)}},X_{D\cup J}))$-null set and thus
a $P(X_{B_1^{(i)}},X_C|\doit(X_{B_2^{(i)}},X_{D\cup J})) \otimes \mu_{B_2^{(i)}}$-null set.
The absolute continuity:
 \[ \mu_{B \cup C} \ll  P(X_{B_1^{(i)}},X_C|\doit(X_{B_2^{(i)}},X_{D\cup J})) \otimes \mu_{B_2^{(i)}}, \]
 then renders $N^{(i)}$ a $\mu_{B \cup C}$-null set. This shows that the finite union:
 \[ N := \bigcup_{i \in I} N^{(i)}, \]
 is a $\mu_{B \cup C}$-null set as well.  Again, note that on the complement $N^\cmpl$ all Markov kernels agree with $Q(X_A|X_B,X_C,X_D)$ and are thus all equal on $N^\cmpl$. \\

3.) Consider the Markov kernel $Q(X_A|X_C,X_D)$ from Corollary \ref{cor:do-calculus} point 3 and
for $i \in I$ the measurable set:
\begin{align*}
    \tilde N^{(i)} := \bigl\{ x_{B^{(i)} \cup C \cup D \cup J} &\in \Xcal_{B^{(i)} \cup C \cup D \cup J} \,\big|\,
        Q(X_A|X_C=x_C,X_D=x_D)\\ 
       &\neq P(X_A|\doit(X_{B^{(i)}}=x_{B^{(i)}}),X_C=x_C,\doit(X_{D\cup J}=x_{D \cup J})) \bigr\}.
\end{align*}
Again, by the essential uniqueness of conditional Markov kernels, Theorem \ref{thm-regular-conditional-Markov-kernel}, the set $\tilde N^{(i)}$ is a $P(X_C|\doit(X_{B^{(i)}},X_{D\cup J}))$-null set. By the absolut continuity $ \mu_C \ll P(X_C|\doit(X_{B^{(i)}},X_{D\cup J}))$ we get that $\tilde N^{(i)}$ is a $\mu_C$-null set. This shows that the measurable set:
\begin{align*}
    N^{(i)} := \bigl\{ x_{B \cup C \cup D \cup J} &\in \Xcal_{B \cup C \cup D \cup J} \,\big|\,
        Q(X_A|X_C=x_C,X_D=x_D)\\ 
       &\neq P(X_A|\doit(X_{B^{(i)}}=x_{B^{(i)}}),X_C=x_C,\doit(X_{D\cup J}=x_{D \cup J})) \bigr\}.
\end{align*}
is a $\mu_C$-null set as well. Then the finite union:
 \[ N := \bigcup_{i \in I} N^{(i)}, \]
 is also a $\mu_C$-null set. 
Again, note that on the complement $N^\cmpl$ all Markov kernels agree with $Q(X_A|X_C,X_D)$ and are thus all equal on $N^\cmpl$.
\end{proof}
\end{Cor}


We now summarize on how to apply Corollary \ref{cor:as-do-calculus} more concretely in its own theorem.

\begin{Thm}[Almost-sure do-calculus]
    \label{thm:as-do-calculus}
        Consider an IL-BCN:
    \[M= \lp G^+=\lp J,(V,U),E^+ \rp, \lp P_v ( X_v|\doit( X_{\Pa^{G^+}(v)})) \rp_{v \in V \cup U} \rp.\]
        Further assume that we have reference measures $\mu_v$ on $\Xcal_v$ for every $v \in V$. 
    We then put $\mu_F := \bigotimes_{v \in F} \mu_v$ for $F \ins V$.
    Let $A,B,C \ins V$ and $J \ins D \ins J \cup V$ be such that $A,B,C,D$ are pairwise disjoint.
    Then we have the following 3 rules relating marginal conditional to marginal interventional Markov kernels:
    \begin{enumerate}
        \item \underline{Insertion/deletion of observation}:
            Assume that we want to establish the a.s.-equality:
            \[ P(X_A|X_B,X_C,\doit(X_{D})) = P(X_A|X_C,\doit(X_{D})) \qquad \mu_{B \cup C}\text{-a.s.}, \]
         then it is sufficient to assume/check the following d-separation and absolute continuities:
         \begin{align*} 
             A &\Perp^d_{G_{\doit(D)}} B \given C \cup D, &
             \mu_{B \cup C} &\ll P(X_B,X_C|\doit(X_{D}))\ll \mu_{B \cup C}.
         \end{align*}
     \item \underline{Action/observation exchange}: 
         Assume that we want to establish the a.s.-equality:
         \[ P(X_A|X_B,X_C,\doit(X_{D}))  = P(X_A|\doit(X_B),X_C,\doit(X_{D})) 
         \qquad \mu_{B \cup C}\text{-a.s.}, \]
         then it is sufficient to assume/check the following d-separation and absolute continuities:
         \begin{align*} 
             A & \Perp^d_{G_{\doit(I_B,D)}} I_B \given B \cup C \cup D, &          
                    \mu_{B \cup C} &\ll  P(X_B,X_C|\doit(X_{D})) \ll \mu_{B \cup C},\\
                 &&   \mu_C &\ll  P(X_C|\doit(X_B,X_{D})) \ll \mu_C.
         \end{align*}
     \item \underline{Insertion/deletion of action}: 
         Assume that we want to establish the a.s.-equality:
         \[ P(X_A|\doit(X_B),X_C,\doit(X_{D})) = P(X_A|X_C,\doit(X_{D})) \qquad \mu_C\text{-a.s.}, \]
          then it is sufficient to assume/check the following d-separation and absolute continuities:
         \begin{align*} 
             A &\Perp^d_{G_{\doit(I_B,D)}} I_B \given C  \cup D, &          
             \mu_C &\ll P(X_C|\doit(X_B,X_{D})) \ll \mu_C, \\
                   && \mu_C &\ll P(X_C|\doit(X_{D})) \ll \mu_C.
         \end{align*}
    \end{enumerate}
 Note that the two-sided absolute continuities hold if the indicated Markov kernel has a strictly positive Doob-Radon-Nikodym derivative w.r.t.\ the corresponding reference product measure.
 \begin{proof}
     This follows directly from Corollary \ref{cor:as-do-calculus}.
 \end{proof}
\end{Thm}

\begin{Rem}
 Note that in Theorem \ref{thm:as-do-calculus} (in contrast to Corollaries \ref{cor:do-calculus} and \ref{cor:as-do-calculus}) we cannot easily formulate the independence of variables $X_{J \sm D}$ in the presented way. 
 If we accept that then we can simplify the formulas (as done in Theorem \ref{thm:as-do-calculus}) and assume that $J \ins D$ and thus $D \cup J = D$, which makes then the d-separation requirements weaker (due to extra conditioning on $J$).
\end{Rem}

\begin{comment}
\begin{Thm}[Almost-sure do-calculus]
    \label{thm:as-do-calculus}
        Consider an IL-BCN:
    \[M= \lp G^+=\lp J,(V,U),E^+ \rp, \lp P_v ( X_v|\doit( X_{\Pa^{G^+}(v)})) \rp_{v \in V \cup U} \rp.\]
        Further assume that we have reference measures $\mu_v$ on $\Xcal_v$ for every $v \in V$. 
    We then put $\mu_F := \bigotimes_{v \in F} \mu_v$ for $F \ins V$.
    Let $A,B,C \ins V$ and $D \ins J \cup V$ be such that $A,B,C,D$ are pairwise disjoint.
    Then we have the following 3 rules relating marginal conditional to marginal interventional Markov kernels:
    \begin{enumerate}
        \item \underline{Insertion/deletion of observation}:
            Assume that we want to establish the a.s.-equality:
            \[ P(X_A|X_B,X_C,\doit(X_{D\cup J})) = P(X_A|X_C,\doit(X_{D\cup J})) \qquad \mu_{B \cup C}\text{-a.s.}, \]
         then it is sufficient to assume/check the following d-separation and absolute continuities:
         \begin{align*} 
             A &\Perp^d_{G_{\doit(D)}} B \given C \cup D, &
             \mu_{B \cup C} &\ll P(X_B,X_C|\doit(X_{D\cup J}))\ll \mu_{B \cup C}.
         \end{align*}
     \item \underline{Action/observation exchange}: 
         Assume that we want to establish the a.s.-equality:
         \[ P(X_A|X_B,X_C,\doit(X_{D\cup J}))  = P(X_A|\doit(X_B),X_C,\doit(X_{D\cup J})) 
         \qquad \mu_{B \cup C}\text{-a.s.}, \]
         then it is sufficient to assume/check the following d-separation and absolute continuities:
         \begin{align*} 
             A & \Perp^d_{G_{\doit(I_B,D)}} I_B \given B \cup C \cup D, &          
                    \mu_{B \cup C} &\ll  P(X_B,X_C|\doit(X_{D\cup J})) \ll \mu_{B \cup C},\\
                 &&   \mu_C &\ll  P(X_C|\doit(X_B,X_{D\cup J})) \ll \mu_C.
         \end{align*}
     \item \underline{Insertion/deletion of action}: 
         Assume that we want to establish the a.s.-equality:
         \[ P(X_A|\doit(X_B),X_C,\doit(X_{D\cup J})) = P(X_A|X_C,\doit(X_{D\cup J})) \qquad \mu_C\text{-a.s.}, \]
          then it is sufficient to assume/check the following d-separation and absolute continuities:
         \begin{align*} 
             A &\Perp^d_{G_{\doit(I_B,D)}} I_B \given C  \cup D, &          
             \mu_C &\ll P(X_C|\doit(X_B,X_{D\cup J})) \ll \mu_C, \\
                   && \mu_C &\ll P(X_C|\doit(X_{D\cup J})) \ll \mu_C.
         \end{align*}
    \end{enumerate}
 Note that the two-sided absolute continuities hold if the indicated Markov kernel has a strictly positive Doob-Radon-Nikodym derivative w.r.t.\ the corresponding reference product measure.
 \begin{proof}
     This follows directly from Corollary \ref{cor:as-do-calculus}.
 \end{proof}
\end{Thm}

\begin{Rem}
Note that in Theorem \ref{thm:as-do-calculus} (in contrast to Corollaries \ref{cor:do-calculus} and \ref{cor:as-do-calculus})  we can w.l.o.g.\ assume that $J \ins D$ and thus $D \cup J = D$, which makes the d-separation requirements weaker (due to extra conditioning on $J$).
\end{Rem}
\end{comment}
